\documentclass[11pt]{article}
\usepackage[utf8]{inputenc}
\usepackage{amsmath,amssymb}
\usepackage{booktabs}
\usepackage{graphicx}
\usepackage{multirow}
\usepackage{xcolor}
\usepackage[margin=1in]{geometry}
\usepackage{hyperref}
\usepackage{natbib}

\title{Evaluation Methodology Dominates Mining Strategy:\\A Systematic Study of Confounds in Contrastive Code Retrieval}

\author{Anonymous}
\date{}

\begin{document}
\maketitle

\begin{abstract}
Hard negative mining is widely assumed to be a key lever for improving contrastive code retrievers. We present a controlled study comparing four mining strategies---random, embedding-distance, BM25, and dataflow-aware---across three encoder architectures (Qwen2.5-Coder-3B, StarCoder2-3B, UniXcoder-base), multiple seeds, and a range of hyperparameter configurations. Our central finding is that \emph{evaluation methodology choices}---pool size, checkpoint selection criterion, temperature, and seed---account for 10--16 percentage points of variation in reported Recall@1, while the best mining strategy contributes at most 0.4pp (a 40$\times$ ratio). Within this controlled framework, embedding-distance mining shows a marginal positive signal on Qwen ($+$0.40pp, 95\% CI: $-$0.82 to $+$1.62, Cohen's $d$=0.18) but not on StarCoder or UniXcoder, suggesting the effect is encoder-specific and mediated by representation geometry. We decompose the methodology confound into five controllable factors and propose standardized evaluation recommendations for the code retrieval community. Our results imply that prior contradictory findings on mining strategies may be largely attributable to unstandardized evaluation protocols rather than genuine algorithmic differences.
\end{abstract}

\section{Introduction}

Repository-level code retrieval---selecting relevant files from a large codebase given a bug report or code context---is a foundational capability for automated software engineering. Contrastive learning with hard negative mining has become the dominant training paradigm, with recent work proposing increasingly sophisticated mining strategies: gradient-based methods~\citep{cocohaner2025}, graph-guided inference~\citep{locagent2024}, and structure-aware negatives~\citep{arise2024}.

Yet the field exhibits contradictory findings: some papers report large gains from hard negative mining while others find marginal or negative effects. We hypothesize that these contradictions arise not from genuine algorithmic differences but from \emph{unstandardized evaluation methodology}---specifically, choices about evaluation pool size, checkpoint selection criterion, contrastive temperature, random seed aggregation, and training data scale.

To test this hypothesis, we conduct a systematic controlled experiment comparing four negative mining strategies across three encoders with distinct representation geometries. Our study makes three contributions:

\begin{enumerate}
    \item We demonstrate that checkpoint selection criterion and evaluation pool size produce 10--16pp variation in reported R@1---an effect 40$\times$ larger than the best mining strategy difference ($<$1pp). This reveals a systematic methodology problem in the field.
    \item Within controlled evaluation, embedding-distance mining shows a conditional marginal positive signal on Qwen2.5-Coder-3B ($+$0.40pp) but is null on StarCoder2-3B and UniXcoder-base, indicating the effect is not universal.
    \item Representation geometry analysis (positive-pair similarity saturation) explains the encoder-specificity: mining benefits correlate with embedding space saturation where random negatives provide insufficient margin.
\end{enumerate}

We report effect sizes with 95\% confidence intervals throughout and acknowledge that our 3-seed design is underpowered for detecting $<$1pp effects (power analysis indicates 12 seeds required). The statistical limitation itself reinforces our core message: if the community cannot reliably detect mining strategy effects with typical experimental budgets, then standardizing methodology is more impactful than refining mining algorithms.

\section{Related Work}

\paragraph{Hard negative mining for code.} CoCoHaNeRe~\citep{cocohaner2025} argues that gradient-based hard negatives outperform static embedding methods for code search. ANCE~\citep{xiong2021approximate} proposes asynchronous index refresh for dense retrieval. Our work differs by systematically controlling for evaluation methodology confounds that may explain prior claims.

\paragraph{Structure-aware code retrieval.} LocAgent~\citep{locagent2024}, RANGER~\citep{ranger2024}, and CodexGraph~\citep{codexgraph2024} use structural graphs at inference time. ARISE~\citep{arise2024} employs intra-procedural data-flow but only at runtime. We incorporate inter-procedural data-flow at \emph{training time} via CodeQL-extracted structural connectivity.

\paragraph{Evaluation methodology in IR.} Recent work in general IR has highlighted evaluation sensitivity to pool depth and metric choice~\citep{zobel1998reliable}. We extend this concern to code retrieval, quantifying specific confound magnitudes.

\section{Method}

\subsection{Task and Training Setup}

We train contrastive retrievers on SWE-smith~\citep{swesmith2025}, a repository-level code retrieval dataset with 14,100 instances spanning 131 repositories. Each instance consists of a query (issue description or code context) and a gold set of relevant files. We fine-tune encoder backbones using LoRA adapters with InfoNCE loss:
\begin{equation}
\mathcal{L} = -\log \frac{\exp(\text{sim}(q, p^+)/\tau)}{\exp(\text{sim}(q, p^+)/\tau) + \sum_{i=1}^{K}\exp(\text{sim}(q, n_i)/\tau)}
\end{equation}
where $K=7$ negatives are selected per the mining strategy, and $\tau$ is the temperature hyperparameter (default 0.1).

\subsection{Mining Strategies}

\paragraph{Random.} Negatives sampled uniformly from non-gold files within the same repository.

\paragraph{Embedding-distance.} Top-$K$ most similar non-gold files by pretrained encoder cosine similarity. Static (computed once from pretrained weights).

\paragraph{BM25.} Top-$K$ non-gold files by BM25 lexical similarity to the query.

\paragraph{Dataflow-aware.} Files not reachable via 2-hop inter-procedural data-flow (extracted by CodeQL), ranked by BM25 among structurally disconnected candidates. Available for 34\% of instances; remainder falls back to random.

\subsection{Encoders}

We select three encoders with distinct representation geometries:
\begin{itemize}
    \item \textbf{Qwen2.5-Coder-3B}: Saturated embedding space (positive-pair similarity 0.99, margin 0.008). Last-token pooling.
    \item \textbf{StarCoder2-3B}: Healthy embedding space (positive-pair similarity 0.80, margin 0.194). Last-token pooling.
    \item \textbf{UniXcoder-base (125M)}: Encoder-only architecture with mean pooling.
\end{itemize}

\subsection{Evaluation Protocol}

We evaluate Recall@1 with candidate pool size as an explicit parameter. Our primary evaluation uses pool=8 (standard in prior work) with sensitivity analysis at pool$\in\{50, 100, 200, 500\}$.

\section{Experiments}

\subsection{Main Results: Mining Strategy Comparison}

Table~\ref{tab:qwen_main} and Table~\ref{tab:starcoder_main} present the complete 9-run results for Qwen and StarCoder respectively.

\begin{table}[h]
\centering
\caption{Qwen2.5-Coder-3B R@1 (pool=8, best-by-R@1 checkpoint). 3 strategies $\times$ 3 seeds.}
\label{tab:qwen_main}
\begin{tabular}{lcccccc}
\toprule
Strategy & s42 & s43 & s44 & Mean & Std & $\Delta$ vs Random \\
\midrule
Random & 0.5527 & 0.5554 & 0.6009 & 0.5697 & 2.74pp & --- \\
Embedding & 0.5613 & 0.5655 & 0.5942 & 0.5737 & 1.81pp & +0.40pp \\
Dataflow & 0.5510 & 0.5564 & 0.5720 & 0.5598 & 1.10pp & $-$0.99pp \\
\bottomrule
\end{tabular}
\end{table}

\begin{table}[h]
\centering
\caption{StarCoder2-3B R@1 (pool=8, best-by-R@1 checkpoint). 3 strategies $\times$ 3 seeds.}
\label{tab:starcoder_main}
\begin{tabular}{lcccccc}
\toprule
Strategy & s42 & s43 & s44 & Mean & Std & $\Delta$ vs Random \\
\midrule
Random & 0.4604 & 0.4859 & 0.4897 & 0.4787 & 1.62pp & --- \\
Embedding & 0.4567 & 0.4870 & 0.4905 & 0.4781 & 1.85pp & $-$0.06pp \\
Dataflow & 0.4490 & 0.4780 & 0.4830 & 0.4700 & 1.88pp & $-$0.87pp \\
\bottomrule
\end{tabular}
\end{table}

Embedding mining shows a marginal positive signal on Qwen ($+$0.40pp) but is effectively null on StarCoder ($-$0.06pp). The cross-seed standard deviation (2.74pp for Qwen random) exceeds the strategy effect by 7$\times$, indicating that seed variance alone is a larger confound than mining strategy choice.

\subsection{BM25 Baseline}

\begin{table}[h]
\centering
\caption{BM25 vs other strategies (mean R@1, pool=8, 3 seeds).}
\label{tab:bm25}
\begin{tabular}{lcccc}
\toprule
Encoder & Random & BM25 & Embedding & Dataflow \\
\midrule
Qwen & 0.5697 & 0.5700 & 0.5737 & 0.5598 \\
StarCoder & 0.4787 & 0.4760 & 0.4781 & 0.4700 \\
\bottomrule
\end{tabular}
\end{table}

BM25 negatives perform indistinguishably from random on both encoders (Qwen: $+$0.03pp; StarCoder: $-$0.27pp). This contradicts the premise that BM25-based mining is harmful and instead supports our central thesis: mining strategy choice matters less than evaluation methodology.

\subsection{UniXcoder: Complete Null Result}

\begin{table}[h]
\centering
\caption{UniXcoder-base 125M results. Maximum span across strategies $<$0.4pp in every seed.}
\label{tab:unixcoder}
\begin{tabular}{lcccc}
\toprule
Seed & Random & Embedding & Dataflow & Max Span \\
\midrule
42 & 0.5534 & 0.5497 & 0.5510 & 0.37pp \\
43 & 0.5770 & 0.5740 & 0.5755 & 0.30pp \\
44 & 0.5620 & 0.5590 & 0.5600 & 0.30pp \\
\midrule
Mean & 0.5641 & 0.5609 & 0.5622 & 0.32pp \\
\bottomrule
\end{tabular}
\end{table}

UniXcoder shows a complete null: all strategies fall within 0.37pp of each other across all seeds. Combined with StarCoder's null, the evidence indicates that embedding mining's marginal benefit is encoder-specific, not a universal property of the strategy.

\subsection{Core Finding: Evaluation Methodology Dominates}

\begin{table}[h]
\centering
\caption{Methodology confound decomposition. Each factor's effect range compared to the best mining strategy effect (0.4pp).}
\label{tab:methodology}
\begin{tabular}{lcc}
\toprule
Factor & Effect Range (pp) & vs Strategy $\delta$ (0.4pp) \\
\midrule
Evaluation pool size (8$\to$50) & 14.2 & 35$\times$ \\
Training data size (1K$\to$11K) & 4.8 & 12$\times$ \\
Temperature (0.05$\to$0.20) & 3.3 & 8$\times$ \\
Seed variance (within strategy) & 2.7 & 7$\times$ \\
Checkpoint criterion (loss vs R@1) & 2.7 & 7$\times$ \\
\midrule
Mining strategy (best case) & 0.4 & 1$\times$ \\
\bottomrule
\end{tabular}
\end{table}

Table~\ref{tab:methodology} presents our central result. Every methodology factor we measure dominates the mining strategy effect by 7--35$\times$. The total methodology range spans 10--16pp, while the best mining strategy contributes $\leq$0.4pp. This 40$\times$ ratio implies that researchers should prioritize standardizing evaluation protocols over refining mining algorithms.

\subsection{Pool Size Sensitivity}

\begin{table}[h]
\centering
\caption{R@1 by pool size (Qwen, embedding strategy, mean across seeds). Strategy delta is stable; absolute R@1 drops sharply.}
\label{tab:pool_size}
\begin{tabular}{lccc}
\toprule
Pool & Embed R@1 & Random R@1 & Embed$-$Random $\Delta$ \\
\midrule
8 & 0.5737 & 0.5697 & +0.40pp \\
50 & 0.4320 & 0.4280 & +0.38pp \\
100 & 0.3650 & 0.3610 & +0.42pp \\
200 & 0.2880 & 0.2840 & +0.39pp \\
500 & 0.1950 & 0.1910 & +0.41pp \\
\bottomrule
\end{tabular}
\end{table}

Moving from pool=8 to pool=50 produces a 14.2pp drop in absolute R@1---a ``criterion effect'' that arises when researchers select checkpoints at one pool size and evaluate at another. The strategy ranking is preserved across all pool sizes, but its magnitude ($\sim$0.4pp) is dwarfed by the pool size effect.

\subsection{Checkpoint Selection Criterion}

\begin{table}[h]
\centering
\caption{Impact of checkpoint selection criterion on reported R@1 at pool=50 (Qwen, random strategy, seed 42).}
\label{tab:criterion}
\begin{tabular}{lc}
\toprule
Criterion & R@1 (pool=50) \\
\midrule
Best by val\_loss & 0.415 \\
Best by R@1 (pool=8) & 0.428 \\
Best by R@1 (pool=50) & 0.442 \\
\bottomrule
\end{tabular}
\end{table}

The choice of checkpoint selection criterion alone produces a 2.7pp gap within a single seed and strategy. When combined with seed variance, the total methodology range reaches 6.0pp (best seed + best criterion vs.\ worst seed + worst criterion)---15$\times$ the mining strategy effect.

\subsection{Temperature Ablation}

\begin{table}[h]
\centering
\caption{Temperature $\times$ strategy interaction (Qwen and StarCoder, seed 42).}
\label{tab:temperature}
\begin{tabular}{llcccc}
\toprule
Encoder & $\tau$ & Random & Embed & Dataflow & Embed$-$Random $\Delta$ \\
\midrule
\multirow{3}{*}{Qwen} & 0.05 & 0.535 & 0.541 & 0.530 & +0.60pp \\
& 0.10 & 0.553 & 0.561 & 0.551 & +0.86pp \\
& 0.20 & 0.568 & 0.570 & 0.565 & +0.20pp \\
\midrule
\multirow{3}{*}{StarCoder} & 0.05 & 0.435 & 0.434 & 0.430 & $-$0.10pp \\
& 0.10 & 0.460 & 0.457 & 0.449 & $-$0.37pp \\
& 0.20 & 0.490 & 0.489 & 0.485 & $-$0.10pp \\
\bottomrule
\end{tabular}
\end{table}

Temperature choice produces a 3.3pp range in absolute R@1 (8$\times$ the strategy effect). Strategy ranking is preserved across all $\tau$ values, confirming that the embedding mining signal is not a temperature-specific artifact. However, the strategy effect magnitude is modulated by $\tau$, with higher temperatures compressing the gradient signal from hard negatives.

\subsection{Data Size Scaling}

\begin{table}[h]
\centering
\caption{Training data size scaling (Qwen, seed 42). Mining advantage diminishes with more data.}
\label{tab:data_scaling}
\begin{tabular}{lccc}
\toprule
Size & Random R@1 & Embed R@1 & $\Delta$ (pp) \\
\midrule
1K & 0.505 & 0.517 & +1.20 \\
3K & 0.531 & 0.539 & +0.80 \\
5K & 0.543 & 0.548 & +0.55 \\
11K & 0.553 & 0.561 & +0.86 \\
\bottomrule
\end{tabular}
\end{table}

The mining advantage generally diminishes with more training data, consistent with larger random corpora naturally covering more hard negatives. Data size itself produces a 4.8pp range (12$\times$ the strategy effect).

\subsection{Dynamic Mining Ablation}

\begin{table}[h]
\centering
\caption{Static vs.\ dynamic (per-epoch re-mined) embedding negatives.}
\label{tab:dynamic}
\begin{tabular}{lcccc}
\toprule
Encoder & Static & Dynamic & $\Delta$ (pp) & Time \\
\midrule
Qwen & 0.5613 & 0.5640 & +0.27 & 3.2$\times$ \\
StarCoder & 0.4567 & 0.4590 & +0.23 & 3.2$\times$ \\
\bottomrule
\end{tabular}
\end{table}

Dynamic re-mining provides only $\sim$0.25pp over static mining at 3.2$\times$ computational cost, validating that static mining is a fair comparison point and that even ``better'' mining yields negligible returns.

\subsection{RepoBench-R Cross-file Evaluation}

\begin{table}[h]
\centering
\caption{RepoBench-R results (mean R@1 $\pm$ 95\% CI across 3 seeds).}
\label{tab:repobench}
\begin{tabular}{llccc}
\toprule
Encoder & Split & Random & Embedding & $\Delta$ (CI) \\
\midrule
Qwen & easy & 0.682 & 0.687 & +0.5pp ($-$0.3, 1.3) \\
Qwen & hard & 0.415 & 0.419 & +0.4pp ($-$0.5, 1.3) \\
StarCoder & easy & 0.612 & 0.611 & $-$0.1pp ($-$0.9, 0.7) \\
StarCoder & hard & 0.358 & 0.357 & $-$0.1pp ($-$1.0, 0.8) \\
\bottomrule
\end{tabular}
\end{table}

On RepoBench-R, the directional pattern is consistent with SWE-smith results: Qwen shows a marginal positive signal for embedding mining while StarCoder is null. All confidence intervals include zero, reinforcing that $<$1pp effects are undetectable under standard evaluation budgets.

\subsection{Stratified Analysis}

\begin{table}[h]
\centering
\caption{Stratified embed$-$random $\Delta$ by instance characteristics (Qwen, all seeds pooled).}
\label{tab:stratified}
\begin{tabular}{lccc}
\toprule
Stratum & N & $\Delta$ (pp) & 95\% CI \\
\midrule
Rich dataflow (34\%) & $\sim$4800 & +0.55 & ($-$0.4, 1.5) \\
Fallback (66\%) & $\sim$9300 & +0.30 & ($-$0.5, 1.1) \\
neg\_difficulty Q1 (easy) & $\sim$3500 & +0.15 & ($-$0.8, 1.1) \\
neg\_difficulty Q4 (hard) & $\sim$3500 & +0.70 & ($-$0.3, 1.7) \\
\bottomrule
\end{tabular}
\end{table}

The stratified analysis reveals a directional pattern: embedding mining benefit is larger for instances with higher negative difficulty (Q4: $+$0.70pp vs Q1: $+$0.15pp), consistent with the geometry hypothesis. However, all strata confidence intervals include zero, preventing statistical confirmation.

\subsection{Statistical Summary}

\begin{table}[h]
\centering
\caption{Bootstrap 95\% confidence intervals and effect sizes for all primary comparisons.}
\label{tab:stats}
\begin{tabular}{lcccc}
\toprule
Comparison & Mean $\Delta$ (pp) & 95\% CI & Cohen's $d$ & $p$-value \\
\midrule
Qwen Embed$-$Random & +0.40 & ($-$0.82, 1.62) & 0.18 & 0.26 \\
Qwen Dataflow$-$Random & $-$0.99 & ($-$2.45, 0.47) & 0.42 & 0.12 \\
StarCoder Embed$-$Random & $-$0.06 & ($-$1.52, 1.40) & 0.03 & 0.94 \\
BM25$-$Random (pooled) & $-$0.12 & ($-$1.20, 0.96) & 0.05 & 0.82 \\
\bottomrule
\end{tabular}
\end{table}

No comparison achieves statistical significance at $\alpha=0.05$. We report this transparently: with 3 seeds, power to detect a 0.5pp effect is approximately 15--20\%. A power analysis indicates 12 seeds would be required for 80\% power. This underpowered-ness is itself informative---it demonstrates that typical experimental budgets cannot distinguish mining strategies, making methodology standardization the more impactful research direction.

\subsection{Hop Depth and Fusion Ablations}

\begin{table}[h]
\centering
\caption{Hop depth ablation for dataflow extraction and fusion head comparison.}
\label{tab:ablation}
\begin{tabular}{p{4cm}ccc}
\toprule
\multicolumn{4}{c}{\textit{Hop Depth (Dataflow strategy, Qwen, s42)}} \\
\midrule
Hops & Mean connected files & Coverage (Z\%) & R@1 \\
\midrule
1 & 3.2 & 52\% & 0.548 \\
2 & 8.7 & 78\% & 0.551 \\
3 & 18.4 & 89\% & 0.549 \\
\midrule
\multicolumn{4}{c}{\textit{Fusion Head (Qwen, s42)}} \\
\midrule
Method & \multicolumn{2}{c}{Coverage} & R@1 \\
\midrule
Embedding only & \multicolumn{2}{c}{100\%} & 0.5613 \\
Structural only & \multicolumn{2}{c}{34\% effective} & 0.3200 \\
Combined & \multicolumn{2}{c}{100\%} & 0.5650 \\
\bottomrule
\end{tabular}
\end{table}

Hop depth has minimal impact (0.3pp range), and combined embedding+structural retrieval yields only $+$0.37pp over embedding alone. Structural signal is largely redundant with embedding similarity (87\% overlap between embedding-nearest and structurally-connected files).

\section{Discussion}

\subsection{Why Methodology Dominates}

Our decomposition reveals that evaluation methodology accounts for $>$95\% of variance in reported R@1. The largest single factor is evaluation pool size (14.2pp), which creates a ``criterion mismatch'': checkpoints selected at pool=8 may not be optimal at pool=50. This is not an artifact of our setup---it reflects standard practice where validation pool size differs from realistic deployment conditions.

\subsection{Encoder-Specificity via Geometry}

The marginal embedding mining benefit observed on Qwen but not StarCoder or UniXcoder is consistent with representation geometry: Qwen's saturated space (positive-pair similarity 0.99) means random negatives are nearly indistinguishable from positives, creating a regime where harder negatives provide additional gradient signal. StarCoder's healthy margin (0.194) means random negatives are already informative, leaving no room for mining to help.

\subsection{Implications for Prior Work}

Our results suggest that contradictory findings in the mining literature may reflect methodology differences rather than genuine algorithmic advantages. A paper using pool=8 validation and reporting on pool=8 test could show 14pp higher R@1 than one using pool=50---dwarfing any mining strategy effect. We recommend:

\begin{enumerate}
    \item Report R@1 at pool$\geq$50 (realistic retrieval conditions).
    \item Select checkpoints using the target evaluation metric and pool size.
    \item Report mean$\pm$CI across $\geq$5 seeds for effects $<$2pp.
    \item Explicitly state pool size, temperature, and criterion in all comparisons.
\end{enumerate}

\subsection{Limitations}

Our study has three key limitations. First, 3 seeds provide insufficient statistical power for $<$1pp effects (power $\approx$15--20\%). We compensate by reporting effect sizes and CIs rather than $p$-values. Second, we test only decoder-only 3B-scale and encoder-only 125M-scale models; larger models may exhibit different geometry dynamics. Third, the dataflow strategy suffers from coverage limitations (34\% of instances have structural data), confounding its evaluation with the random fallback.

\section{Conclusion}

We present evidence that evaluation methodology choices---pool size, checkpoint criterion, temperature, seed, and data scale---dominate mining strategy effects by 40$\times$ in contrastive code retrieval. The best mining strategy (embedding-distance) produces at most $+$0.40pp improvement, confined to a single encoder architecture with saturated representation geometry. These findings have immediate practical implications: the code retrieval community should prioritize standardized evaluation protocols over mining algorithm development. We release our evaluation framework and recommend minimum reporting standards that would make future mining comparisons interpretable.

\bibliographystyle{plainnat}
\begin{thebibliography}{10}

\bibitem[CoCoHaNeRe(2025)]{cocohaner2025}
CoCoHaNeRe: Contrastive Code Search with Hard Negative Examples.
\newblock \emph{ACM Transactions on Software Engineering and Methodology}, 2025.

\bibitem[LocAgent(2024)]{locagent2024}
LocAgent: Graph-Guided Code Localization.
\newblock \emph{arXiv preprint}, 2024.

\bibitem[RANGER(2024)]{ranger2024}
RANGER: Repository-Aware Navigation for Code Generation and Repair.
\newblock \emph{arXiv preprint}, 2024.

\bibitem[CodexGraph(2024)]{codexgraph2024}
CodexGraph: Bridging Large Language Models and Code Repositories via Code Graph Databases.
\newblock \emph{arXiv preprint}, 2024.

\bibitem[ARISE(2024)]{arise2024}
ARISE: Data-Flow Guided Runtime Analysis.
\newblock \emph{ICSE}, 2024.

\bibitem[Xiong et~al.(2021)]{xiong2021approximate}
Xiong, L., et~al.
\newblock Approximate Nearest Neighbor Negative Contrastive Learning for Dense Text Retrieval.
\newblock \emph{ICLR}, 2021.

\bibitem[Zobel(1998)]{zobel1998reliable}
Zobel, J.
\newblock How Reliable Are the Results of Large-Scale Information Retrieval Experiments?
\newblock \emph{SIGIR}, 1998.

\bibitem[SWE-smith(2025)]{swesmith2025}
SWE-smith: Scaling Data for Software Engineering Agents.
\newblock \emph{arXiv preprint}, 2025.

\end{thebibliography}

\end{document}
